\documentclass{article}
\usepackage[T1]{fontenc}
\usepackage{lmodern}
\usepackage{graphicx}
\usepackage{amsmath,amssymb}
\usepackage[a4paper,margin=2cm]{geometry}
\usepackage[absolute,overlay]{textpos}
\setlength{\TPHorizModule}{1mm}
\setlength{\TPVertModule}{1mm}
\usepackage[indonesian]{babel}
\usepackage{hyperref}
\setlength{\emergencystretch}{2em}
% unit: Halaman 4

\begin{document}

% === Halaman 4 (llm) ===

\begin{itemize}
  \item Oleh karena itu, dikenal dua macam \textit{cipher} di dalam kriptografi klasik:
  \begin{enumerate}
    \item \textit{Cipher} Substitusi (\textit{substitution Cipher})
    \begin{itemize}
      \item metode enkripsi dan dekripsi menggunakan teknik substitusi
    \end{itemize}
    \item \textit{Cipher} Transposisi (\textit{transposition Cipher})
    \begin{itemize}
      \item metode enkripsi dan dekripsi menggunakan teknik transposisi
    \end{itemize}
  \end{enumerate}
  \item Kombinasi kedua teknik tersebut membentuk \textit{product cipher} atau \textit{super enkripsi}
  \[ 
    product\ cipher = cipher\ substitusi + cipher\ transposisi
  \]
\end{itemize}

\end{document}
